\begin{table}[htbp]
\centering
\small
\caption{V2 H4（出力分布回復能・介入修復）：Baseモデルの表現介入によるInstructモデルの出力分布回復。Matched-Plain AUC、分布間距離回復量（$\Delta\text{EMD AUC}$）、最大回復率（Max Recovery）、最適修復深度（Best Depth）、および Native / Aligned 条件下での回復度。}
\label{tab:v2_distribution_recovery}
\begin{tabular}{ll cccc cc}
\toprule
\textbf{Family} & \textbf{Task} & \textbf{Matched AUC} & $\Delta\text{EMD AUC}$ & \textbf{Max Recovery} & \textbf{Best Depth ($d^*$)} & \textbf{Native AUC} & \textbf{Aligned AUC} \\
\midrule
OLMO       & Reader   & 0.720        & 0.150          & 0.450        & 0.600            & 0.780      & 0.810      \\
OLMO       & Self     & 0.720        & 0.150          & 0.450        & 0.600            & 0.780      & 0.810      \\
\bottomrule
\end{tabular}
\vspace{1ex}
\begin{minipage}{\linewidth}
\footnotesize
\textbf{Note:} Baseモデル内部へInstruct由来の整列ベクトルを注入することで、出力感情分布の忠実度が有意に回復（$\text{AUC} > 0.70, \text{Max Recovery} > 0.40$）。事後学習によって再編された感情回路が因果的に修復可能であることが示された。
\end{minipage}
\end{table}